\makeatletter
% Layout macros for slide decks that started life in PowerPoint.
% Coordinates and sizes are PostScript points (bp) from the top-left
% corner of a 960 x 540 bp page, the size PowerPoint uses for 16:9.
\geometry{papersize={960bp,540bp}}
\usepackage{fontspec}
\usepackage{tikz}
\usepackage{varwidth}
\usepackage{etoolbox}
\setbeamertemplate{navigation symbols}{}
\setbeamertemplate{headline}{}
\setbeamertemplate{footline}{}
\setbeamersize{text margin left=0bp, text margin right=0bp}
\setbeamercolor{background canvas}{bg=white}
\setbeamercolor{normal text}{fg=black}
\hypersetup{colorlinks=false, pdfborder={0 0 0}}
% Org turns a bare address inside a link text into \url; keep the text font.
\urlstyle{same}

% Fonts: each family key used by the decks gets a \f<key> switch.
% A deck picks what it needs in its own style file.
\newcommand{\pptxfamily}[3][]{%
  \IfFontExistsTF{#2}%
    {\expandafter\newfontfamily\csname f#3\endcsname{#2}}%
    {\expandafter\newfontfamily\csname f#3\endcsname{#1}}}

% ---------------------------------------------------------------- text boxes
\newsavebox{\tbbox}
\newcount\tbdepth
\pgfkeys{/tb/.is family, /tb,
  x/.store in=\tbx, y/.store in=\tby, s/.store in=\tbs, l/.store in=\tbl,
  c/.store in=\tbc, f/.store in=\tbf, p/.store in=\tbp, a/.store in=\tba,
  b/.code={\def\tbbold{\bfseries}}, i/.code={\def\tbital{\itshape}},
  v/.code={\def\tbverb{1}}, bb/.code={\def\tbbb{\bfseries}},
  ind/.store in=\tbind, ind2/.store in=\tbindtwo, off2/.store in=\tbofftwo,
  bc/.store in=\tbbc, bf/.store in=\tbbf, bs/.store in=\tbbs, bs2/.store in=\tbbstwo,
  w/.store in=\tbw, r/.store in=\tbr, fw/.store in=\tbfw,
  % code: a code listing; the deck's \pptxcode family keeps quotes and
  % hyphens as typed (no TeX ligatures).
  code/.code={\def\tbf{\pptxcode}},
  defaults/.code={%
    \def\tbx{0}\def\tby{0}\def\tbs{18}\def\tbl{21.6}\def\tbc{black}%
    \def\tbf{\pptxbase}\def\tbp{0}\def\tba{l}\def\tbbold{}\def\tbital{}%
    \def\tbverb{0}\def\tbind{18}\def\tbindtwo{18}\def\tbofftwo{0}%
    \def\tbbc{}\def\tbbf{}\def\tbbs{\textbullet}\def\tbbstwo{}\def\tbw{}\def\tbr{0}\def\tbfw{}\def\tbbb{\mdseries}%
    \providecommand\pptxcode{\pptxbase}}}

\newcommand{\tblabel}[1]{%
  {\ifx\tbbf\empty\else\csname f\tbbf\endcsname\fi
   \ifx\tbbc\empty\else\color{\tbbc}\fi\tbbb\upshape #1}}
\newcommand{\tblists}{%
  \renewenvironment{itemize}{%
    \global\advance\tbdepth by 1
    \ifnum\tbdepth=1
      \def\tbcurlabel{\tblabel{\tbbs}}%
      \begin{list}{\tbcurlabel}{%
        \setlength\leftmargin{\tbind bp}\setlength\labelwidth{\tbind bp}%
        \setlength\labelsep{0pt}\setlength\itemsep{\tbp bp}%
        \setlength\topsep{0pt}\setlength\partopsep{0pt}\setlength\parsep{0pt}%
        \setlength\itemindent{0pt}\setlength\listparindent{0pt}%
        \def\makelabel####1{####1\hfil}}%
    \else
      \ifx\tbbstwo\empty\def\tbcurlabel{\tblabel{\tbbs}}\else\def\tbcurlabel{\tblabel{\tbbstwo}}\fi
      \begin{list}{\tbcurlabel}{%
        \setlength\leftmargin{\dimexpr\tbofftwo bp+\tbindtwo bp\relax}%
        \setlength\labelwidth{\tbindtwo bp}%
        \setlength\labelsep{0pt}\setlength\itemsep{\tbp bp}%
        \setlength\topsep{\tbp bp}\setlength\partopsep{0pt}\setlength\parsep{0pt}%
        \setlength\itemindent{0pt}\setlength\listparindent{0pt}%
        \def\makelabel####1{####1\hfil}}%
    \fi}%
    {\end{list}\global\advance\tbdepth by -1 }%
  % Numbered lists: same geometry, the label is the item number.
  \renewenvironment{enumerate}{%
    \global\advance\tbdepth by 1
    \begin{list}{\tblabel{\arabic{enumi}.}}{\usecounter{enumi}%
      \setlength\leftmargin{\tbind bp}\setlength\labelwidth{\tbind bp}%
      \setlength\labelsep{0pt}\setlength\itemsep{\tbp bp}%
      \setlength\topsep{0pt}\setlength\partopsep{0pt}\setlength\parsep{0pt}%
      \setlength\itemindent{0pt}\setlength\listparindent{0pt}%
      \def\makelabel####1{####1\hfil}}}%
    {\end{list}\global\advance\tbdepth by -1 }}

\NewDocumentEnvironment{tb}{O{}}{%
  \begingroup
  \pgfkeys{/tb/.cd, defaults, #1}%
  \global\tbdepth=0
  \def\tbanchor{base west}%
  \if c\tba \def\tbanchor{base}\fi
  \if r\tba \def\tbanchor{base east}\fi
  \begin{lrbox}{\tbbox}%
  % Font and colour go outside the [t] box: a colour change at the top of
  % its vertical list would move the box baseline above the first line.
  \csname f\tbf\endcsname
  \fontsize{\tbs bp}{\tbl bp}\selectfont\tbbold\tbital\color{\tbc}%
  \ifnum\tbverb=1 \begin{minipage}[t]{900bp}\else\begin{varwidth}[t]{940bp}\fi
  \setlength\parskip{\tbp bp}\setlength\parindent{0pt}%
  \if c\tba\centering\fi \if r\tba\raggedleft\fi \if l\tba\raggedright\fi
  \tblists
  \ifnum\tbverb=1 \def\verbatim@font{\normalfont\ttfamily\fontsize{\tbs bp}{\tbl bp}\selectfont}\fi
}{%
  \ifnum\tbverb=1 \end{minipage}\else\end{varwidth}\fi
  \end{lrbox}%
  % fw: squeeze or stretch the text to the width it had in the PDF
  \ifx\tbfw\empty\else
    \sbox{\tbbox}{\resizebox{\tbfw bp}{\height}{\usebox{\tbbox}}}%
  \fi
  \begin{tikzpicture}[remember picture, overlay]
    \node[inner sep=0pt, outer sep=0pt, anchor=\tbanchor, rotate=\tbr]
      at ([shift={(\tbx bp,-\tby bp)}]current page.north west) {\usebox{\tbbox}};
  \end{tikzpicture}%
  \endgroup\ignorespacesafterend}

% Shapes and pictures go on the background canvas, under the title.
% The frame body runs before the page is shipped out, so the file name
% is kept in a global macro that the background template reads.
\gdef\slidegfx{}
\newcommand{\slidegraphics}[1]{\gdef\slidegfx{#1}}
% \pptximg{width}{height}{file}: xdvipdfmx embeds every \includegraphics
% anew, so a picture that recurs on many slides would be stored many
% times.  This draws a file the first time and reuses that copy after.
\newcount\pptx@imgn
\newcommand\pptximg[3]{%
  \ifcsname pptx@img@#3\endcsname
    \def\pptx@do{uxobj}\def\pptx@file{}%
  \else
    \global\advance\pptx@imgn 1
    \expandafter\xdef\csname pptx@img@#3\endcsname{\the\pptx@imgn}%
    \def\pptx@do{image}\def\pptx@file{(#3)}%
  \fi
  \vbox to #2{\vss\hbox to #1{%
    \special{pdf:\pptx@do\space @pptximg\csname pptx@img@#3\endcsname\space
      width #1 height #2 \pptx@file}\hss}}}
\setbeamertemplate{background}{\ifx\slidegfx\empty\else\input{\slidegfx}\fi}
\BeforeBeginEnvironment{frame}{\gdef\slidegfx{}}

% ---------------------------------------------------------------- titles
% Frame options: tx, ty (baseline), ts (size), tc (colour), tf (family),
% tw (wrap width), tl (line pitch), notitle, pg (page in the old PDF).
\define@key{beamerframe}{pg}{}
\define@key{beamerframe}{tx}{\def\ftx{#1}}
\define@key{beamerframe}{ty}{\def\fty{#1}}
\define@key{beamerframe}{ts}{\def\fts{#1}}
\define@key{beamerframe}{tc}{\def\ftc{#1}}
\define@key{beamerframe}{tf}{\def\ftf{#1}}
\define@key{beamerframe}{tw}{\def\ftw{#1}}
\define@key{beamerframe}{tl}{\def\ftl{#1}}
\define@key{beamerframe}{notitle}[]{\def\fnotitle{1}}
\newcommand{\titledefaults}{%
  \def\ftx{\pptxtitlex}\def\fty{\pptxtitley}\def\fts{\pptxtitles}%
  \def\ftc{black}\def\ftf{\pptxtitlef}\def\ftw{900}\def\ftl{0}\def\fnotitle{0}}
\titledefaults
\BeforeBeginEnvironment{frame}{\titledefaults}
\setbeamertemplate{frametitle}{%
  \ifnum\fnotitle=0
  \begin{tikzpicture}[remember picture, overlay]
    \node[inner sep=0pt, anchor=base west, text=\ftc]
      at ([shift={(\ftx bp,-\fty bp)}]current page.north west)
      {\begin{varwidth}[t]{\ftw bp}\raggedright
        \csname f\ftf\endcsname\bfseries
        \ifdim\ftl pt=0pt \fontsize{\fts bp}{1.15\dimexpr\fts bp\relax}\else\fontsize{\fts bp}{\ftl bp}\fi
        \selectfont\insertframetitle\end{varwidth}};
  \end{tikzpicture}%
  \fi}
\makeatother
